% !TEX root = chapter-standalone.tex

\section{Preorders}

\label{sec:tradeoffs-ordered-sets}
\linkvideo{spring2021-tradeoffs:tradeoffs:orders:pre-pos-tot} % Pre-orders, partial orders, and total orders

%\linkvideo{spring2021-tradeoffs:tradeoffs:orders} % Orders

%
% While we discussed how categories can describe the way in which one resource can be turned into another, this kind of modelling did not allow for quantitative statements.
% For example, it is good to know that we can obtain motion from electric power, but, how fast can we go with a certain amount of power?

% To achieve a quantitative theory, we need to specify various degrees of resources and functionality.
% One way of doing this is using the idea of orders.

% Such orderings arise naturally in engineering as criteria for judging whether one design is better or worse than another.
% As an example, suppose you need to buy a battery.
% In this simple example, you can think of having two resources: mass and money.
% A lighter battery might be more expensive, and a heavier one could be more affordable.
% How to choose among batteries, if you do not prefer one resource over the other?
% How to model this?
% We introduce these concepts by adding levels of specificity.


\todotext{Discuss resource convertibility to motivate preorders}

\begin{marginfigure}
    \centering
    \includesag{pre_order_graph}
    \caption[A pre-order represented as a graph.]{A \SY{pre-order} represented as a graph.}
    \label{fig:pre_order_graph}
\end{marginfigure}

A \SY{pre-order} is a set together with a binary relation that is both \SY{reflexive} (\cref{def:endo_reflexive_irreflexive}) and \SY{transitive}~(\cref{def:transitive-relation}).
\begin{ctdefinition}[Pre-ordered set]
    \label{def:preorder}
    A \maindef{pre-ordered set} is a tuple~$\posAdefinition$, where~$\posAset$ is a set, called the \emph{carrier set} or \emph{underlying set}, together with a relation~$\posAleq$ that is
    \SY{reflexive} and \SY{transitive}.
\end{ctdefinition}
An example of a \SY{pre-ordered set} represented as a graph is shown in \cref{fig:pre_order_graph}.
In the graph representation of a \SY{pre-order}~\posA, we draw an arrow between~$\posela$ and~$\poselb$ if~$\posela \posAleq \poselb$.

\begin{example}
    The reachability relationship in any directed graph (potentially including cycles) is a pre-order.
    The pre-order $\posA$ is defined as follows.
    The set $\posAset$ is the set of nodes of the graph.
    Take any two nodes~$\ela,\elb \setin \posAset$.
    One has $\ela \posAleq \elb$ if and only if there is a path from $\ela$ to $\elb$ in the directed graph.
    There is always a path from a node to itself (reflexivity), and given a path from $\ela$ to $\elb$, and one from $\elb$ to $\elc$, we know that there is a path from $\ela$ to $\elc$ (transitivity).
\end{example}

\begin{exercise}
    \label{ex:findpreorders}
    Consider the set $\posAset=\makeset{\ela,\elb,\elc}$.
    Which of the following are pre-orders? Why?
    \begin{enumerate}
        \item $\posA= \makeset{\tup{\ela,\ela},\tup{\ela,\elb},\tup{\elb,\ela},\tup{\elb,\elb}}$.
        \item $\posB= \makeset{\tup{\ela,\elb},\tup{\elb,\elc},\tup{\elc,\ela}}$.
    \end{enumerate}
\end{exercise}

\begin{solution}
    $\posA$ is a pre-order, because it satisfies reflexivity and transitivity.
    $\posB$ violates both reflexivity (e.g., $\tup{\ela,\ela}$ is missing), and transitivity (e.g., $\tup{\ela,\elc}$ is missing).
\end{solution}

\todotext{JL: we should include more examples here. }


